\documentclass[11pt,a4paper]{article}

\usepackage[T1]{fontenc}
\usepackage[utf8]{inputenc}
\usepackage[spanish,es-tabla]{babel}
\usepackage{amsmath}
\usepackage[margin=2.4cm]{geometry}
\usepackage{booktabs}
\usepackage{array}
\usepackage{longtable}
\usepackage{enumitem}
\usepackage{parskip}
\usepackage{caption}
\usepackage{xcolor}
\usepackage{tcolorbox}
\usepackage{fancyhdr}
\usepackage[hidelinks]{hyperref}

\newcolumntype{L}[1]{>{\raggedright\arraybackslash}p{#1}}

% Atajo para nombres de archivo / codigo con guion bajo
\newcommand{\code}[1]{\texttt{#1}}
\newcommand{\fpath}[1]{\texttt{#1}}

% Caja de "ojo / nota"
\newtcolorbox{nota}{colback=yellow!8,colframe=yellow!55!black,
  boxrule=0.5pt,arc=2pt,left=6pt,right=6pt,top=4pt,bottom=4pt}
\newtcolorbox{regla}{colback=red!5,colframe=red!55!black,
  boxrule=0.6pt,arc=2pt,left=6pt,right=6pt,top=4pt,bottom=4pt}

\pagestyle{fancy}
\fancyhf{}
\rhead{\small Manual del tablero de research}
\lhead{\small Uso personal}
\cfoot{\thepage}
\renewcommand{\headrulewidth}{0.3pt}

\title{\textbf{Tablero de research personal}\\[3pt]
\large Manual de uso --- de la puesta en marcha a cada funcionalidad}
\author{Herramienta personal de apoyo a la decisi\'on}
\date{Junio de 2026}

\begin{document}
\maketitle
\thispagestyle{fancy}

\noindent\rule{\textwidth}{0.4pt}

\small
\textbf{Qu\'e es esto.} Un manual de la herramienta de research que corro localmente
(Python + Streamlit, base SQLite). Cubre c\'omo ponerla en funcionamiento, qu\'e hace
cada solapa de la app y para qu\'e sirve cada cosa adentro, y tambi\'en el trabajo que
queda \emph{fuera} de la app (informes, prompts, daily digest, score de la cartera y
los archivos de datos). \emph{No es asesoramiento financiero}: las decisiones de
inversi\'on las tomo yo; la herramienta solo organiza datos y ayuda a razonar ideas.
\normalsize

\noindent\rule{\textwidth}{0.4pt}

\tableofcontents
\vspace{0.4cm}

%==============================================================
\section{Mapa general}

La herramienta tiene dos mitades que conviene tener claras desde el principio:

\begin{enumerate}[leftmargin=1.4em]
  \item \textbf{La app (el tablero).} Un programa en Python con interfaz web que se
  abre en el navegador. Tiene seis solapas: \emph{Scorecard}, \emph{Screener},
  \emph{Precio}, \emph{Bonos}, \emph{Informes} y \emph{Carteras simuladas}. Baja
  precios y ratios sola, calcula la TIR de los bonos, rankea acciones y sigue
  carteras simuladas.
  \item \textbf{El trabajo fuera de la app (en el chat y por terminal).} Ac\'a entra
  el an\'alisis de prensa y de informes macro (que produce los informes diarios), los
  prompts reutilizables, el \emph{daily digest} de las carteras, el \emph{score} de
  una cartera puntual, y el mantenimiento de los archivos de datos (los CSV).
\end{enumerate}

\noindent\textbf{Piezas de software (todas en la misma carpeta):}

\begin{longtable}{L{4.4cm} L{9.9cm}}
\toprule
\textbf{Archivo} & \textbf{Para qu\'e sirve} \\
\midrule
\endhead
\fpath{app.py} & La interfaz (Streamlit). Es lo que se ejecuta para abrir el tablero. \\
\fpath{portfolio\_research.py} & El motor de datos: base SQLite, descarga de precios y ratios, c\'alculo de TIR/duration/paridad de bonos, scorecard, carteras simuladas, informes. \\
\fpath{screener\_score.py} & El motor del \emph{score compuesto} 0--10 del screener (6 dimensiones, pesos por estilo y horizonte). \\
\fpath{score\_cartera.py} & Script de terminal: puntu\'a las acciones de una cartera con el score compuesto. \\
\fpath{daily\_digest.py} & Script de terminal: resumen diario de las carteras (retorno, alpha, se\~nales). \\
\fpath{seed\_cartera\_*.py} & Scripts que ``siembran'' carteras simuladas predefinidas. \\
\fpath{universe\_cedears.csv} & Universo de acciones/CEDEARs ($\sim$200 tickers). \\
\fpath{ratios\_cedears.csv} & Ratio de conversi\'on CEDEAR por acci\'on (para precios en pesos). \\
\fpath{universe\_bonds.csv} & Universo de bonos (24: soberanos, BOPREAL y ONs). \\
\fpath{cashflows.csv} & Flujo de fondos de cada bono (para calcular TIR/duration/paridad). \\
\fpath{research.db} & La base SQLite. Se crea sola la primera vez; guarda todo lo descargado. \\
\fpath{informes/} & Carpeta con los informes diarios (\fpath{.md}, opcional \fpath{.pdf}/\fpath{.tex}). \\
\bottomrule
\end{longtable}

\begin{regla}
\textbf{Regla de oro del proyecto:} \emph{nunca inventar datos financieros} (flujos,
tickers, TIR). Si falta el flujo de un bono nuevo, se pide la captura del Cashflow +
Datos t\'ecnicos de Balanz y se valida la TIR calculada contra la de Balanz antes de
confiar. Los scripts de carteras tampoco inventan precios: toman el \'ultimo cierre
cargado y, si falta, saltean y avisan.
\end{regla}

%==============================================================
\section{Puesta en marcha}

\subsection{Requisitos}

Corre todo en Windows, en VS Code, dentro de la carpeta del proyecto:

\begin{quote}
\fpath{C:\textbackslash\allowbreak Users\textbackslash\allowbreak Luca\textbackslash\allowbreak Desktop\textbackslash\allowbreak INVERSIONES\textbackslash\allowbreak CODIGO\textbackslash}
\end{quote}

\noindent Hace falta Python instalado y las librer\'ias:

\begin{quote}
\code{pip install streamlit yfinance pandas numpy}
\end{quote}

\noindent (Lo ideal es hacerlo dentro de un entorno virtual: \code{python -m venv venv}
y activarlo, para no mezclar dependencias. Una vez creado, alcanza con activarlo cada
vez antes de correr la app.)

\subsection{Abrir el tablero}

Desde la terminal, parada en la carpeta del proyecto:

\begin{quote}
\code{streamlit run app.py}
\end{quote}

Esto abre el tablero en el navegador (t\'ipicamente \code{http://localhost:8501}). La
base \fpath{research.db} se crea sola la primera vez (la app llama a \code{init\_db()},
\code{seed\_recommendations()} y \code{seed\_bonds()} al arrancar). Para cerrar, se
corta el proceso en la terminal (Ctrl+C).

\subsection{Primera vez: cargar datos}

La base arranca vac\'ia de precios. En la barra de arriba del tablero:

\begin{enumerate}[leftmargin=1.4em]
  \item Apretar \textbf{``Actualizar precios (universo)''}. Baja el \'ultimo a\~no de
  cierres de todo el universo, en lotes (de a 40 tickers). Tarda un rato.
  \item Traer los ratios en tandas: elegir cu\'antos por tanda (campo \emph{``Ratios
  por tanda''}, por defecto 50) y apretar \textbf{``Traer N ratios''}. Es
  \emph{reanudable}: se puede cortar y seguir despu\'es; saltea los que ya baj\'o hoy.
  Repetir hasta llegar a 200/200 (el contador de cobertura est\'a arriba).
\end{enumerate}

\subsection{Alternativa por terminal}

El motor tambi\'en se maneja sin abrir la app, \'util para automatizar:

\begin{longtable}{L{6.2cm} L{8.1cm}}
\toprule
\textbf{Comando} & \textbf{Qu\'e hace} \\
\midrule
\endhead
\code{python portfolio\_research.py update} & Baja precios de todo el universo. \\
\code{python portfolio\_research.py funds} & Baja ratios (reanudable). \\
\code{python portfolio\_research.py track} & Imprime el scorecard en la terminal. \\
\code{python portfolio\_research.py all} & Las tres cosas seguidas. \\
\bottomrule
\end{longtable}

%==============================================================
\section{La barra superior (com\'un a toda la app)}

Arriba de las solapas, siempre visible, est\'a el panel de control de datos:

\begin{itemize}[leftmargin=1.4em]
  \item \textbf{Cobertura:} ``Universo acciones: N $\cdot$ Ratios cargados: X/N''.
  Cu\'antos tickers hay y de cu\'antos ya se bajaron ratios.
  \item \textbf{\rmfamily ``Actualizar precios (universo)'':} baja un a\~no de cierres
  ajustados de todos los tickers (en lotes de 40). Llena la tabla \code{prices}.
  \item \textbf{``Ratios por tanda'' + ``Traer N ratios'':} baja los ratios
  fundamentales de yfinance de a tandas. Reanudable.
\end{itemize}

\begin{nota}
Los precios y ratios vienen de \textbf{yfinance} (Yahoo Finance). Es la base de las
solapas Scorecard, Screener, Precio y Carteras simuladas. Si una solapa muestra
``todav\'ia no hay precios/ratios'', es porque falta correr estos botones.
\end{nota}

%==============================================================
\section{Solapa 1 --- Scorecard}

\textbf{Para qu\'e sirve:} medir, de forma objetiva, c\'omo vienen las recomendaciones
de analistas que cargo (por ejemplo de Allaria), compar\'andolas contra el mercado
(SPY) desde la fecha de cada recomendaci\'on.

Cada fila es una recomendaci\'on; las columnas:

\begin{longtable}{L{2.4cm} L{11.7cm}}
\toprule
\textbf{Columna} & \textbf{Significado} \\
\midrule
\endhead
Fuente & De qui\'en es la recomendaci\'on (ej.\ ``Allaria CEDEARS''). \\
Fecha & Fecha de la recomendaci\'on (la ventana de medici\'on arranca ah\'i). \\
Ticker & La acci\'on recomendada. \\
PrecioRec & Precio al momento de la recomendaci\'on. \\
Actual & \'Ultimo precio cargado en la base. \\
Ret\% & Retorno desde la recomendaci\'on: $(\text{Actual}/\text{PrecioRec}-1)\times100$. \\
Bench\% & Lo que hizo el SPY en la misma ventana (desde la fecha de la recomendaci\'on). \\
Alpha\% & Ret\% $-$ Bench\%. Cu\'anto le gan\'o (o perdi\'o) al mercado. \\
aObj\% & Avance al precio objetivo: qu\'e porcentaje del camino al \emph{target} ya recorri\'o, $\frac{\text{Actual}-\text{PrecioRec}}{\text{Objetivo}-\text{PrecioRec}}\times100$ (s\'olo si la recomendaci\'on ten\'ia objetivo). \\
\bottomrule
\end{longtable}

\noindent\textbf{C\'omo se cargan recomendaciones nuevas.} Hoy vienen ``sembradas'' en
una lista dentro de \fpath{portfolio\_research.py} (\code{SEED\_RECOMMENDATIONS}), que
se inserta en la tabla \code{recommendations} al arrancar. Para sumar una nueva, se
agrega una fila a esa lista con el formato

\begin{quote}
\code{(fuente, fecha, ticker, acci\'on, precio\_al\_llamado, precio\_objetivo)}
\end{quote}

\noindent y se reinicia la app. El \code{precio\_objetivo}
puede ir vac\'io si no hay target.

%==============================================================
\section{Solapa 2 --- Screener}

\textbf{Para qu\'e sirve:} filtrar y rankear el universo de $\sim$200 CEDEARs por
ratios, con un \emph{score compuesto} que resume la calidad de cada empresa en un
n\'umero. Es la herramienta para pasar de ``200 acciones'' a una lista corta.

\subsection{El bloque ``Score compuesto''}

Arriba de todo (expander abierto por defecto) se elige con qu\'e \emph{lente} se
puntu\'a el universo:

\begin{itemize}[leftmargin=1.4em]
  \item \textbf{Estilo:} \emph{quality, value, GARP, growth, dividend, defensive}.
  Cambia qu\'e dimensiones pesan m\'as (un perfil \emph{value} prioriza valuaci\'on
  barata; \emph{growth}, crecimiento y momentum; etc.).
  \item \textbf{Horizonte:} \emph{long / medium / short-term}. Multiplica los pesos:
  en \emph{short-term} el momentum pesa mucho m\'as y la calidad menos; en
  \emph{long-term}, al rev\'es.
  \item \textbf{Modo:} \emph{Por sector} (recomendado) o \emph{Umbrales absolutos}.
  Por sector, cada empresa se compara con las de su rubro (un banco con bancos); los
  sectores con menos de 5 empresas caen a umbrales fijos.
\end{itemize}

El score se calcula sobre \emph{todo} el universo (para que los percentiles por sector
usen todos los pares) y agrega la columna de \emph{momentum}. El detalle de c\'omo se
arma est\'a en la secci\'on~\ref{sec:score}.

\subsection{Canastas tem\'aticas}

Botones de un toque (del informe del 23/06) que filtran el universo a un grupo de
tickers con los ratios y el orden ya seteados. Por ejemplo ``IA semis \& net'',
``Megacap de-rating'', ``Bancos AR'', ``O\&G AR'', etc. Los tickers de la canasta que
no tengan ratios cargados no aparecen (y se avisan abajo). Se quitan con
``\,Quitar canasta\,''. Para editarlas o agregar nuevas se toca el diccionario
\code{THEMATIC\_BASKETS} en \fpath{app.py}.

\subsection{Presets y filtros manuales}

\begin{itemize}[leftmargin=1.4em]
  \item \textbf{Presets} (botones): \emph{Value} (P/E y P/BV bajos), \emph{Calidad}
  (ROE y margen altos), \emph{Dividendos} (yield m\'inimo) y \emph{Limpiar} (saca
  todos los filtros). Setean los checks y valores; despu\'es se pueden ajustar a mano.
  \item \textbf{Filtros manuales:} P/E m\'aximo, P/BV m\'aximo, FV/EBITDA m\'aximo,
  ROE\% m\'inimo, Margen\% m\'inimo, Div.Yield m\'inimo (en crudo, ej.\ 0,03 = 3\%).
  Cada uno se activa con su check.
  \item \textbf{Sector:} multiselecci\'on para quedarse solo con ciertos rubros.
  \item \textbf{Ordenar por:} Score o cualquier columna num\'erica, ascendente o
  descendente.
\end{itemize}

La tabla de abajo muestra cu\'antas empresas pasan los filtros y todas las columnas de
ratios m\'as el score y sus subcomponentes. Las columnas num\'ericas son: P/E, P/E
fwd, P/BV, FV/EBITDA, ROE\%, Margen\%, Div.Yield, MktCap, CrecVtas\%, CrecGan\%,
Deuda/Eq, Liq.Corr y Mom\% (momentum).

\subsection{Qu\'e mide cada columna (en criollo)}

Esta es la gu\'ia de las columnas del screener, pensada para alguien que no viene de
finanzas. La idea de fondo: cada n\'umero contesta una pregunta distinta sobre la
empresa. Hay tres familias: \emph{cu\'an cara est\'a} (valuaci\'on), \emph{qu\'e tan
buena es} (rentabilidad y crecimiento), \emph{qu\'e tan s\'olida es} (solidez), y
\emph{c\'omo se viene moviendo} (momentum). Ninguna sirve sola; se leen en conjunto.

\subsubsection*{Valuaci\'on: \textquestiondown{}est\'a cara o barata?}

\begin{itemize}[leftmargin=1.4em]
  \item \textbf{P/E (precio / ganancias).} Cu\'antos a\~nos de las ganancias
  \emph{actuales} de la empresa pag\'as al comprar la acci\'on. Un P/E de 20 significa
  que pag\'as 20 veces lo que la empresa gana en un a\~no. \emph{Por qu\'e lo usamos:}
  es la vara m\'as r\'apida de ``caro vs.\ barato''. Un P/E bajo puede ser una ganga o
  una empresa con problemas; uno alto suele reflejar que el mercado espera mucho
  crecimiento (o que est\'a cara). Solo tiene sentido entre empresas que ganan plata
  (si la ganancia es negativa, no aplica).
  \item \textbf{P/E fwd (P/E \emph{forward}, a futuro).} Lo mismo, pero usando las
  ganancias \emph{estimadas} del pr\'oximo a\~no en lugar de las pasadas. \emph{Por
  qu\'e lo usamos:} si el P/E fwd es m\'as bajo que el P/E normal, el mercado espera
  que las ganancias \emph{crezcan} (la misma acci\'on se vuelve ``m\'as barata'' a
  futuro). Es \'util para no castigar a una empresa que hoy gana poco pero va en
  subida.
  \item \textbf{P/BV (precio / valor libros).} Compara el precio de la acci\'on contra
  el \emph{patrimonio neto contable} de la empresa (lo que valdr\'ia en los papeles si
  vendiera todo y pagara sus deudas). Un P/BV de 1 significa que cotiza justo a su
  valor contable; menos de 1, por debajo. \emph{Por qu\'e lo usamos:} es especialmente
  \'util en bancos y empresas con muchos activos tangibles, donde ``los libros''
  significan algo concreto. Un P/BV muy bajo puede ser oportunidad o se\~nal de que el
  mercado desconf\'ia de esos activos.
  \item \textbf{FV/EBITDA (valor de la empresa / ganancia operativa).} El ``FV'' (valor
  de la firma) es lo que costar\'ia comprar la empresa entera, contando tambi\'en su
  deuda; el EBITDA es la ganancia operativa antes de intereses, impuestos y
  amortizaciones (una medida ``cruda'' de cu\'anto genera el negocio). \emph{Por qu\'e
  lo usamos:} a diferencia del P/E, este ratio \emph{incluye la deuda}, as\'i que
  permite comparar de forma m\'as justa empresas con distinto nivel de endeudamiento.
  Menor = relativamente m\'as barata. Es de los m\'as usados por analistas para comparar
  entre s\'i.
\end{itemize}

\subsubsection*{Rentabilidad y crecimiento: \textquestiondown{}qu\'e tan buena es?}

\begin{itemize}[leftmargin=1.4em]
  \item \textbf{ROE\% (rentabilidad sobre el patrimonio).} De cada \$100 que pusieron
  los accionistas, cu\'antos genera la empresa de ganancia al a\~no. Un ROE de 20\%
  quiere decir \$20 por cada \$100. \emph{Por qu\'e lo usamos:} mide qu\'e tan bien la
  empresa usa el capital de sus due\~nos. Un ROE alto y \emph{sostenido} es se\~nal de
  un buen negocio (marca, ventaja competitiva). Ojo: un ROE alt\'isimo a veces viene de
  mucha deuda, no de calidad; por eso se lee junto a Deuda/Eq.
  \item \textbf{Margen\% (margen de ganancia neta).} De cada \$100 que la empresa
  \emph{vende}, cu\'antos le quedan de ganancia limpia despu\'es de todos los costos.
  \emph{Por qu\'e lo usamos:} mide la eficiencia y el ``poder de fijar precios''. Un
  margen alto suele indicar que la empresa puede cobrar caro sin que se le escapen los
  clientes (poca competencia o producto diferenciado). Var\'ia mucho por industria: un
  supermercado tiene margen fino y un software, margen ancho.
  \item \textbf{CrecVtas\% (crecimiento de ventas).} Cu\'anto crecieron los ingresos
  respecto del a\~no anterior. \emph{Por qu\'e lo usamos:} las ventas son la l\'inea
  ``de arriba'' del negocio; sin crecimiento de ventas, es dif\'icil que crezcan las
  ganancias de forma sana. Es el term\'ometro de si el negocio se est\'a expandiendo o
  estancando.
  \item \textbf{CrecGan\% (crecimiento de ganancias).} Lo mismo, pero en la l\'inea ``de
  abajo'': cu\'anto crecieron las \emph{ganancias}. \emph{Por qu\'e lo usamos:} al final
  lo que mueve el precio de una acci\'on en el largo plazo son las ganancias. Que
  crezcan m\'as r\'apido que las ventas indica que la empresa adem\'as se est\'a
  volviendo m\'as eficiente. Puede ser irregular a\~no a a\~no (de ah\'i que se mire
  junto al crecimiento de ventas).
\end{itemize}

\subsubsection*{Solidez: \textquestiondown{}qu\'e tan s\'olida y segura es?}

\begin{itemize}[leftmargin=1.4em]
  \item \textbf{Deuda/Eq (deuda sobre patrimonio).} Cu\'anta deuda tiene la empresa por
  cada \$1 de capital propio. \emph{Por qu\'e lo usamos:} la deuda potencia las
  ganancias en las buenas, pero hunde a la empresa en las malas. Un valor alto = m\'as
  riesgo financiero (m\'as sensible a subas de tasas o a una ca\'ida de ventas). Menor
  es m\'as conservador. Var\'ia por sector (los bancos manejan otra l\'ogica).
  \item \textbf{Liq.Corr (liquidez corriente).} Cu\'antas veces los activos de corto
  plazo (caja, cobranzas, inventario) cubren las deudas de corto plazo. \emph{Por qu\'e
  lo usamos:} responde ``\textquestiondown{}puede pagar lo que vence este a\~no sin ahogarse?''.
  Por encima de 1 puede cubrir sus pasivos inmediatos; por debajo de 1 podr\'ia tener
  apreturas de caja. Es un chequeo b\'asico de salud financiera de corto plazo.
  \item \textbf{MktCap (capitalizaci\'on de mercado).} El valor total de la empresa en
  bolsa: precio de la acci\'on $\times$ cantidad de acciones. \emph{Por qu\'e lo
  usamos:} es el ``tama\~no'' de la empresa. Las grandes (\emph{large caps}) suelen ser
  m\'as estables y predecibles; las chicas (\emph{small caps}), m\'as vol\'atiles, con
  m\'as potencial de crecimiento pero m\'as riesgo. En el score, el tama\~no se usa como
  respaldo de la dimensi\'on de solidez cuando faltan los datos de deuda y liquidez.
\end{itemize}

\subsubsection*{Renta y tendencia}

\begin{itemize}[leftmargin=1.4em]
  \item \textbf{Div.Yield (rendimiento por dividendos).} El dividendo anual que paga la
  empresa, como porcentaje del precio de la acci\'on. Un yield de 4\% significa que, si
  comprs hoy, cobrar\'as $\sim$\$4 al a\~no por cada \$100 invertidos (adem\'as de lo que
  suba o baje la acci\'on). \emph{Por qu\'e lo usamos:} es la renta ``en efectivo'' del
  due\~no. Importa sobre todo para un perfil de dividendos. Un yield muy alto puede ser
  atractivo o una se\~nal de alerta (a veces sube porque la acci\'on se desplom\'o).
  \item \textbf{Mom\% (momentum, retorno 12--1 meses).} Cu\'anto subi\'o o baj\'o la
  acci\'on en el \'ultimo a\~no, salteando el \'ultimo mes. \emph{Por qu\'e lo usamos:}
  hist\'oricamente, lo que viene subiendo tiende a seguir subiendo un tiempo (y
  viceversa); es el factor ``tendencia''. Se saltea el \'ultimo mes a prop\'osito,
  porque en plazos muy cortos suele haber rebotes en contra. Pesa mucho en horizontes
  cortos y poco en los largos (ver el score compuesto).
\end{itemize}

\begin{nota}
Ninguna columna alcanza por s\'i sola: un P/E bajo \emph{con} ROE alto y poca deuda es
muy distinto de un P/E bajo \emph{porque} la empresa est\'a en problemas. El score
compuesto de la pr\'oxima secci\'on existe justamente para combinarlas en un solo
n\'umero comparable.
\end{nota}

\subsection{El score compuesto en detalle}\label{sec:score}

Vive en \fpath{screener\_score.py}. Resume seis dimensiones en un puntaje 0--10
(en la pr\'actica 1--10) y un \emph{Rating} cualitativo. La idea: pesos base por estilo,
multiplicados por horizonte, y \emph{renormalizados fila por fila} seg\'un qu\'e datos
hay (si a una empresa le falta un ratio, esa dimensi\'on queda N/A y los pesos se
reparten entre las que s\'i tiene; nunca se inventa el dato).

\noindent\textbf{Las seis dimensiones y de d\'onde salen:}

\begin{longtable}{L{2.6cm} L{11.5cm}}
\toprule
\textbf{Dimensi\'on} & \textbf{Ratios que usa} \\
\midrule
\endhead
Quality (calidad) & ROE\% y Margen\% (mayor es mejor). \\
Valuation (valuaci\'on) & P/E (forward si est\'a, si no trailing), P/BV y FV/EBITDA (menor es mejor). \\
Income (renta) & Div.Yield. \\
Growth (crecimiento) & CrecVtas\% y CrecGan\%; si faltan, proxy con P/E fwd vs trailing. \\
Safety (solidez) & Deuda/Eq (menor mejor) y Liq.Corr (mayor mejor); si faltan, tama\~no (MktCap). \\
Momentum & Retorno 12--1 meses (columna externa, ver abajo). \\
\bottomrule
\end{longtable}

\noindent\textbf{Modo ``por sector''.} Cada ratio se convierte en un puntaje 1--10
seg\'un su \emph{percentil dentro del sector} (si el sector tiene al menos 5 empresas
v\'alidas). As\'i un P/E de 18 puede ser ``caro'' entre bancos y ``barato'' entre
tecnol\'ogicas. Si el sector es chico o falta, se usan \emph{umbrales fijos} de
respaldo.

\noindent\textbf{Rating.} A partir del Score: \emph{Fuerte} ($\geq 7{,}5$),
\emph{Bueno} ($\geq 5{,}5$), \emph{Regular} ($\geq 3{,}5$), \emph{Flojo} ($<3{,}5$).
Si la cobertura de datos es baja ($<$ la mitad de las dimensiones) agrega
``(poca data)''. La columna \emph{Cobertura} muestra cu\'antas dimensiones se pudieron
calcular (ej.\ 5/6).

\noindent\textbf{Momentum (12--1).} Se calcula desde los precios: retorno entre el
precio de $\sim$12 meses atr\'as y el de $\sim$1 mes atr\'as (se saltea el \'ultimo
mes, est\'andar para evitar la reversi\'on de corto plazo). Un ticker necesita al
menos $\sim$100 ruedas de historia; si no, queda N/A.

\begin{nota}
El score es un \textbf{ranking heur\'istico sobre datos p\'ublicos}, no un veredicto.
Sirve para ordenar y comparar dentro del universo, no para decidir solo.
\end{nota}

%==============================================================
\section{Solapa 3 --- Precio}

\textbf{Para qu\'e sirve:} ver el gr\'afico hist\'orico de una acci\'on del universo, en
d\'olares o estimando el precio del CEDEAR en pesos.

\begin{itemize}[leftmargin=1.4em]
  \item Se elige un ticker (de los que tienen precios descargados) y se ve el cierre
  ajustado del \'ultimo a\~no.
  \item \textbf{Modo USD:} precio de la acci\'on subyacente. Si hay ratio cargado,
  muestra el ratio del CEDEAR (ej.\ ``20:1'').
  \item \textbf{Modo ARS (CEDEAR estimado):} convierte a pesos con la f\'ormula
  \[
    \text{precio CEDEAR} \approx \frac{\text{precio acci\'on (USD)} \times \text{CCL}}{\text{ratio}}
  \]
  donde \emph{ratio} = cu\'antos CEDEARs equivalen a una acci\'on (de
  \fpath{ratios\_cedears.csv}) y el CCL se baja al toque de fuentes p\'ublicas gratis
  (con bot\'on ``Traer CCL'').
\end{itemize}

\begin{nota}
La conversi\'on a pesos usa el \textbf{CCL de hoy} para toda la serie, as\'i que los
valores de fechas pasadas son orientativos (el CCL de cada d\'ia era distinto).
Conviene comparar el \'ultimo valor estimado contra el precio de Balanz, no el
hist\'orico. Si falta el ratio del ticker (puede ser un ADR que se compra local, no
CEDEAR), muestra en USD y avisa.
\end{nota}

%==============================================================
\section{Solapa 4 --- Bonos}

\textbf{Para qu\'e sirve:} seguir 24 bonos argentinos con precio, TIR, paridad y
duration calculados por la herramienta, m\'as la curva TIR-vs-duration y el riesgo
pa\'is. El universo: 11 soberanos (6 Globales ley NY + 5 Bonares ley local),
6 BOPREAL (ley BCRA) y 7 Obligaciones Negociables / ONs (ley NY).

\subsection{C\'omo se actualiza}

El bot\'on \textbf{``Actualizar bonos (precio + TIR)''} hace dos cosas:

\begin{enumerate}[leftmargin=1.4em]
  \item Baja los precios de \textbf{data912.com} (API p\'ublica gratis, delay
  $\sim$2\,h), de los paneles \code{arg\_bonds}, \code{arg\_corp} (ONs) y
  \code{arg\_notes}.
  \item Calcula \textbf{TIR, duration y paridad} de cada bono que tenga flujo cargado,
  usando ese precio m\'as el flujo de fondos.
\end{enumerate}

La leyenda de cobertura indica cu\'antos bonos tienen flujo cargado (\emph{X/24}).

\subsection{La tabla y los gr\'aficos}

\begin{itemize}[leftmargin=1.4em]
  \item \textbf{Tabla:} Bono, Nombre, Ley, Vence, Precio USD, TIR\%, Paridad\%,
  Duration y la hora de la \'ultima actualizaci\'on.
  \item \textbf{Panel completo de data912} (expander): todos los s\'imbolos crudos que
  devuelve la API, por si hace falta debuguear un precio.
  \item \textbf{Curva TIR vs duration:} cada punto es un bono. Los que quedan
  \emph{por encima} de la curva ofrecen m\'as TIR para igual plazo (los relativamente
  baratos).
  \item \textbf{Riesgo pa\'is (EMBI):} se baja autom\'aticamente de ArgentinaDatos
  (gratis, sin clave) con su propio bot\'on; se grafica la serie en puntos b\'asicos.
\end{itemize}

\subsection{El motor de bonos (c\'omo calcula la TIR)}

Esta es la parte clave: la TIR/paridad/duration \emph{no se bajan de ning\'un lado},
las calcula la herramienta. La fuente externa s\'olo aporta el \emph{precio}.

\begin{itemize}[leftmargin=1.4em]
  \item \textbf{TIR:} se busca, por bisecci\'on, la tasa que hace que el valor presente
  de los flujos futuros iguale al precio. Es la tasa interna de retorno del bono al
  precio de hoy. Se descartan resultados disparatados (fuera del rango
  $-5\%$ a $60\%$).
  \item \textbf{Duration:} se reporta la \emph{duration modificada} (Macaulay $/(1+y)$),
  en a\~nos. Mide la sensibilidad del precio a cambios de tasa.
  \item \textbf{Paridad:} precio dividido el \emph{valor t\'ecnico} ($\times 100$). El
  valor t\'ecnico es el capital residual m\'as los intereses corridos del cup\'on en
  curso (convenci\'on 30/360). Una paridad $<100$ indica que cotiza bajo la par.
  \item \textbf{Limpieza de precios:} data912 a veces devuelve una cotizaci\'on en
  pesos en lugar de en d\'olares; el motor descarta los precios implausibles para no
  ensuciar la tabla ni la curva. El precio de data912 ya viene ``por 100 de nominal
  original'', la misma base que los flujos, as\'i que se usa crudo.
\end{itemize}

\begin{regla}
Para sumar un bono nuevo hace falta su \textbf{flujo de fondos} en
\fpath{cashflows.csv}. Se saca de las capturas de la pesta\~na Cashflow de Balanz y se
\textbf{valida}: la TIR que calcula la herramienta tiene que coincidir con la de
Balanz (siempre dio dentro de 1--3 puntos b\'asicos). Reci\'en ah\'i se conf\'ia. El
flujo de cada bono se carga una sola vez.
\end{regla}

%==============================================================
\section{Solapa 5 --- Informes}

\textbf{Para qu\'e sirve:} tener la bit\'acora de research en un solo lugar, dentro del
tablero. Es el archivo de los informes diarios que se generan en el chat.

\begin{itemize}[leftmargin=1.4em]
  \item Los informes viven como archivos \fpath{.md} en la carpeta \fpath{informes/}
  al lado de la app. El formato principal es Markdown, que se lee \emph{adentro} del
  tablero (se renderiza). La convenci\'on de nombre es
  \fpath{informe\_research\_AAAAMMDD.md}.
  \item Si al lado de un \fpath{.md} hay un \fpath{.pdf} o \fpath{.tex} con el
  \emph{mismo nombre}, aparecen como botones de descarga.
  \item \textbf{Subir un informe:} hay un uploader para \fpath{.md}, \fpath{.pdf} o
  \fpath{.tex}; se guardan en la carpeta autom\'aticamente.
  \item \textbf{Buscador:} busca un t\'ermino (ej.\ ``GD41'', ``Vaca Muerta'') en el
  texto de todos los informes y muestra en cu\'ales aparece, con un fragmento.
  \item \textbf{Selector + lectura:} se elige un informe (el m\'as nuevo primero) y se
  lee renderizado; abajo est\'an los botones de descarga disponibles.
\end{itemize}

%==============================================================
\section{Solapa 6 --- Carteras simuladas}

\textbf{Para qu\'e sirve:} simular carteras de tesis (estilo MarketWatch pero local y
sin login) y seguir su evoluci\'on. Sirve para poner a prueba las convicciones de los
informes sin plata real.

\subsection{Gesti\'on de carteras (panel izquierdo)}

\begin{itemize}[leftmargin=1.4em]
  \item \textbf{Nueva cartera:} nombre, etiqueta opcional (ej.\ ``IA / macro''),
  fecha de inicio y notas (la tesis resumida).
  \item \textbf{Listado:} se selecciona una cartera para ver su detalle a la derecha.
  \item \textbf{Editar / eliminar:} cambiar los metadatos o borrar la cartera (con
  confirmaci\'on; al borrarla se borran sus posiciones).
\end{itemize}

\subsection{Detalle de la cartera (panel derecho)}

\textbf{Agregar posici\'on.} Se elige un ticker (de los que tienen precios
descargados), una fecha de entrada, y se define la posici\'on de una de dos formas:

\begin{itemize}[leftmargin=1.4em]
  \item \textbf{Por cantidad de acciones} (\emph{shares}), o
  \item \textbf{Por monto en USD} (la herramienta infiere las acciones impl\'icitas:
  monto $\div$ precio de entrada).
\end{itemize}

El \textbf{precio de entrada} se precarga con el cierre de yfinance de la fecha
elegida (se puede editar, por ejemplo para usar el precio de Balanz). Es el precio del
\emph{subyacente en USD} (la acci\'on, no el CEDEAR).

\textbf{M\'etricas y tabla.} Arriba: cantidad de posiciones, costo base, valor actual y
retorno total. La tabla de posiciones muestra, por cada una: Ticker, fecha de entrada,
shares, costo base, precio de entrada, precio actual, valor en USD, Ret\%, Alpha\%,
Peso\% y la nota. Definiciones:

\begin{itemize}[leftmargin=1.4em]
  \item \textbf{Ret\%:} retorno de la posici\'on desde la entrada.
  \item \textbf{Alpha\%:} Ret\% menos el retorno del SPY en la misma ventana de cada
  posici\'on.
  \item \textbf{Peso\%:} valor de la posici\'on sobre el valor total de la cartera.
\end{itemize}

\textbf{Gr\'afico hist\'orico.} La evoluci\'on del valor de la cartera en USD (semanal
o diario), con el SPY \emph{normalizado} al valor inicial superpuesto, para comparar
la evoluci\'on relativa. Abajo, la distribuci\'on del valor por ticker.

\begin{nota}
Los precios son del subyacente en USD (yfinance), \emph{sin} ajuste de CCL ni ratio
CEDEAR. La cartera no simula bonos: es de acciones/CEDEARs. Es una simulaci\'on de
seguimiento de ideas, no asesoramiento.
\end{nota}

%==============================================================
\section{Trabajo fuera de la app}

Adem\'as del tablero, hay un flujo de an\'alisis que se hace en el chat y por terminal.

\subsection{Los informes: de d\'onde salen}

El an\'alisis se hace en el chat: le paso prensa (FT / WSJ / Bloomberg) o informes
macro (JPM / FMI / Banco Mundial / BCRA) y obtengo \emph{candidatos razonados} dentro
de mi universo. El resultado se vuelca en un informe en Markdown
(\fpath{informe\_research\_AAAAMMDD.md}), se guarda en la carpeta \fpath{informes/} y se
lee desde la solapa Informes. Opcionalmente se exporta a \fpath{.pdf} o \fpath{.tex}.
Los candidatos despu\'es se verifican en el screener (acciones) o en el m\'odulo de
bonos.

\subsection{Los prompts de an\'alisis}

Est\'an en \fpath{prompts\_analisis.md}, listos para pegar en el chat. Son dos:

\begin{enumerate}[leftmargin=1.4em]
  \item \textbf{An\'alisis tem\'atico} (prensa $\rightarrow$ CEDEARs). Toma titulares o
  un art\'iculo y pide: identificar los temas de fondo, mapear la cadena de valor
  (l\'ogica de ``picos y palas''), traducirlo a tickers \emph{con CEDEAR}, dar la tesis
  en 1--2 l\'ineas y qu\'e ratios chequear en el screener, y cerrar con una lista corta
  priorizada y los riesgos en contra.
  \item \textbf{An\'alisis macro argentino} (informe macro $\rightarrow$ bonos +
  acciones). Pide resumir los drivers macro, decir qu\'e implican para el riesgo pa\'is
  (compresi\'on o suba), bajarlo a la curva de bonos (qu\'e tramo se beneficia, por
  duration/convexidad) y a las acciones argentinas, y cerrar con riesgos.
\end{enumerate}

Las reglas de ambos prompts son las mismas que rigen el proyecto: no decir
``comprá'', no inventar tickers ni datos, y ser expl\'icito con la incertidumbre. El
archivo tambi\'en tiene una tabla para el seguimiento peri\'odico de las carteras.

\subsection{El daily digest}

\fpath{daily\_digest.py} imprime, para cada cartera simulada, un resumen del d\'ia:
retorno total, alpha vs SPY, mejor y peor posici\'on, y las \emph{se\~nales} de
disparador que se hayan activado. La idea es correrlo todos los d\'ias y pegar la
salida en el chat para analizarla.

\noindent\textbf{Uso:} primero actualizar precios (en el tablero o
\code{python portfolio\_research.py update}), despu\'es \code{python daily\_digest.py}.

\noindent\textbf{Se\~nales (disparadores mec\'anicos, NO \'ordenes).} Cada posici\'on se
clasifica en un \emph{tier} con su umbral de toma de ganancia (TP) y de revisi\'on de
stop:

\begin{longtable}{L{4.6cm} L{4.2cm} L{4.6cm}}
\toprule
\textbf{Tier} & \textbf{TP / Stop} & \textbf{Qui\'enes} \\
\midrule
\endhead
Core AI-infra & $+25\%$ / $-15\%$ & NVDA, MU, AVGO, AMD, TSM\ldots \\
Especulativos & $+40\%$ / $-25\%$ & OKLO, ALAB, IREN, CRWV, SMCI \\
Contrarian Argentina & $+35\%$ / $-20\%$ & YPF, GGAL, BMA, BBAR, PAM, VIST\ldots \\
\bottomrule
\end{longtable}

\noindent Si una posici\'on cruza su TP o su stop, el digest la marca para
\emph{revisar} (no es una orden de compra/venta; la decisi\'on la tomo yo). Los
umbrales y los tickers de cada tier se editan al principio del script.

\subsection{El score de una cartera}

\fpath{score\_cartera.py} puntu\'a las acciones de una cartera con el mismo score
compuesto del screener. Corre el scoring sobre \emph{todo} el universo (para que los
percentiles por sector tengan sentido) y despu\'es muestra solo los tickers de la
cartera elegida, ordenados por score, con sus subcomponentes (Q/Val/Gro/Saf/Mom) y el
Rating.

\noindent\textbf{Uso:} \code{python score\_cartera.py "Agresiva 3M jun-26"} (sin
argumento, lista las carteras disponibles).

\noindent La lente por defecto es \emph{growth / short-term / por sector}, pensada para
una cartera agresiva de corto plazo (en short-term el momentum pesa m\'as). Si se la
puntuara con quality/long-term, los nombres agresivos se ver\'ian peor a prop\'osito.

\subsection{Sembrar carteras simuladas}

Hay scripts que crean carteras predefinidas de 100.000 USD a partir de las
convicciones de los informes, tomando el precio de entrada del \'ultimo cierre cargado
(no inventan precios; si falta uno, saltean y avisan):

\begin{longtable}{L{5.3cm} L{8.7cm}}
\toprule
\textbf{Script} & \textbf{Cartera que crea} \\
\midrule
\endhead
\fpath{seed\_cartera\_simulada.py} & ``Tesis research jun-26'': diversificada (tem\'atico + macro AR), $\sim$20 posiciones (infra de IA, megacap, rotaci\'on c\'iclica, bloque Argentina, algo de China). \\
\fpath{seed\_cartera\_agresiva.py} & ``Agresiva 3M jun-26'': concentrada en IA-infra, 9 posiciones, m\'as un sleeve contrarian Argentina. \\
\fpath{seed\_cartera\_agresiva\allowbreak\_cp\_.py} & ``Agresiva 3M jun-26(CP)'': versi\'on revisada con el score compuesto (saca OKLO y CRWV por pre-ingresos, suma TSM). \\
\bottomrule
\end{longtable}

\noindent\textbf{Uso:} con precios actualizados, \code{python seed\_cartera\_simulada.py}
(o el que corresponda). Si ya existe una cartera con ese nombre, avisa y no la
duplica; para recrearla hay que borrarla primero desde el tablero.

%==============================================================
\section{Los archivos de datos (CSV)}

Estos cuatro CSV son el ``combustible'' de la herramienta. Se editan a mano para
mantener los universos.

\begin{longtable}{L{4.0cm} L{10.1cm}}
\toprule
\textbf{Archivo} & \textbf{Qu\'e tiene y c\'omo se arma} \\
\midrule
\endhead
\fpath{universe\_cedears.csv} & Una columna \code{ticker}: el universo de
acciones/CEDEARs ($\sim$200). Editarlo agrega o saca acciones del screener, scorecard
y carteras. El benchmark SPY se agrega solo. \\
\addlinespace
\fpath{ratios\_cedears.csv} & \code{ticker, ratio}: cu\'antos CEDEARs equivalen a una
acci\'on. Sirve para estimar el precio en pesos (solapa Precio). Se valida contra
Comafi / BYMA; cambian con los splits. \\
\addlinespace
\fpath{universe\_bonds.csv} & \code{ticker, name, currency, law, maturity}: los 24
bonos (soberanos, BOPREAL, ONs), con su ley y a\~no de vencimiento. \\
\addlinespace
\fpath{cashflows.csv} & \code{ticker, date, coupon, amort} \emph{por 100 de valor
nominal}: el flujo de fondos de cada bono (cup\'on y amortizaci\'on en cada fecha). Es
lo que permite calcular TIR, duration y paridad. Se saca de las capturas de Balanz y
se valida contra su TIR (1--3 pb). Cada bono se carga una sola vez. \\
\bottomrule
\end{longtable}

\noindent\textbf{C\'omo se calculan (resumen de las fuentes):}

\begin{itemize}[leftmargin=1.4em]
  \item \textbf{Precios de acciones/CEDEARs y ratios fundamentales:} autom\'aticos de
  yfinance (botones de la barra superior).
  \item \textbf{Precios de bonos:} autom\'aticos de data912.com (delay $\sim$2\,h).
  \item \textbf{TIR / paridad / duration de bonos:} las calcula la herramienta con el
  flujo (\fpath{cashflows.csv}) m\'as el precio. No se bajan de ning\'un lado.
  \item \textbf{Riesgo pa\'is:} autom\'atico de ArgentinaDatos.
  \item \textbf{CCL (para precios en pesos):} autom\'atico de fuentes p\'ublicas
  (dolarapi / ArgentinaDatos).
\end{itemize}

\begin{nota}
Balanz no tiene API y no se scrapea detr\'as del login. Por eso los flujos de bonos se
cargan a mano desde capturas y se validan. Es la \'unica forma de mantener la regla de
no inventar datos.
\end{nota}

%==============================================================
\section{Rutina y mantenimiento}

\noindent\textbf{D\'ia a d\'ia:}
\begin{enumerate}[leftmargin=1.4em]
  \item Abrir el tablero (\code{streamlit run app.py}) y actualizar precios (y bonos /
  riesgo pa\'is si interesa).
  \item Correr \code{python daily\_digest.py} y pegar la salida en el chat para
  revisar las carteras.
  \item Si hay prensa o informe nuevo, usar el prompt correspondiente, generar el
  informe \fpath{.md}, guardarlo en \fpath{informes/} y verificar candidatos en el
  screener / bonos.
\end{enumerate}

\noindent\textbf{Cada tanto:}
\begin{itemize}[leftmargin=1.4em]
  \item Re-traer ratios para que no queden viejos (botón de tandas).
  \item Sumar un bono: capturar el Cashflow de Balanz, cargarlo en
  \fpath{cashflows.csv}, actualizar bonos y validar la TIR contra Balanz.
  \item Revisar los ratios de CEDEARs si hubo splits.
\end{itemize}

\vspace{0.4cm}
\noindent\rule{\textwidth}{0.4pt}

\small
\textbf{Recordatorio.} Todo lo anterior es una herramienta personal de apoyo a la
decisi\'on. \emph{No es asesoramiento financiero.} Las se\~nales del digest son
disparadores mec\'anicos para revisar, no \'ordenes; el score del screener es un
ranking heur\'istico; y las carteras simuladas son simulaciones. Las decisiones de
inversi\'on las tomo yo.
\normalsize

\end{document}
